\scanpage{216}

\mspage{90}

\uline{Démonstration de} 7.1. On peut donner un énoncé un peu plus délicat, en supposant \inserted{d'emblée} donné, non une partie $M_0$ de $\operatorname{Ob}M$, mais une partie

\[
\mathrm{Mon}_0\subset\mathrm{Mon},\tag*{(7.1)}
\]

satisfaisant la condition suivante

(7.2) Pour \unsure{tte|tt} \deleted{flèche} mono $i\colon A\mathrel{\unsure{-|\to}} B$ dans $M$, il existe \inserted{$i'\colon A'\mathrel{-}B'$} $i'\in\mathrm{Mon}_0$, et une flèche $i\deleted{_0}'\longrightarrow i$ dans \uline{$\mathrm{Fl}$} $M$

\diagram{\begin{tikzcd}[ampersand replacement=\&,row sep=small]
 A' \arrow[r,dash,"i'"] \arrow[d] \& B' \arrow[d]\\
 A \arrow[r,"i"] \& B
\end{tikzcd}}\hfill,

telle que la flèche correspondante \deleted{\unsure{à}}

\[
C=A\amalg_{A'}B'\longrightarrow B
\]

soit un mono, et que $B'\mathrel{\unsure{-|\to}}B$ ne se factorise pas par $A$.

C'est cela qui était visé par la condition (7.1.1) sur $M_0$, qui implique (7.2) sur $\mathrm{Mon}_0$ défini en termes de $M_0$ comme dit au $1^\circ$.

Dans ce qui suit, on suppose donnée une partie $\mathrm{Mon}_0\subset\mathrm{Fl}\,M$ satisfaisant \unsure{(2)|(1)}. En pratique, elle sera souvent déduite d'une partie $M_0$ de $\operatorname{Ob}M$, génératrice pour les monomorphismes, comme \unsure{explicité|expliqué} dans $1^\circ$), mais on n'a pas besoin du fait que $i'\longrightarrow i$ soit un mono (i.e. 7.1.1' est plus fort que ce qu'il nous faut). P.~ex. si $B'\in M_0$ et $B'\longrightarrow B$ est donné, ne se factorisant pas par $A$, on pourra prendre $A'=B'\times_B A$ et $i'\colon A'\longrightarrow B'$ (pour assurer $i'\in\mathrm{Mon}_0$, il faudra \unsure{encore} supposer $M_0$ stable par passage aux sous-objets), et il y a des chances
% archive sheet number: 215
